La trayectoria es
\[
x(t)=e^{-t}\bigl[-0.4\cos(2t)+0.8\sin(2t)\bigr].
\]
El primer máximo positivo ocurre en \(t=\pi/4\) s y vale
\[
x(\pi/4)=0.8e^{-\pi/4}\approx0.364751\ \mathrm{cm}.
\]
Los siguientes extremos alternan de signo y decrecen en magnitud; el menor valor de toda la trayectoria es el inicial, \(-0.40\) cm. Por tanto,
\[
-0.40\le x(t)\le0.8e^{-\pi/4}<0.37,
\]
y se cumple la holgura. La envolvente \(\sqrt{0.8}e^{-t}\) es una cota, no el máximo efectivo.
